% BMS-YOLO: A Lightweight Box-Supervised Pavement Distress Detection Model
% 使用 pdflatex 编译

\documentclass[twocolumn]{sn-jnl}
\usepackage{geometry}
\geometry{left=2.2cm,right=2.2cm,top=2.5cm,bottom=2.5cm}
% 必要宏包
\usepackage{amsmath, amssymb, amsfonts}
\usepackage{graphicx}
\usepackage{epstopdf}
\usepackage{booktabs}
\usepackage{multirow}
\usepackage{array}
\usepackage{hyperref}
\usepackage{xcolor}

% 标题、作者、摘要
\title{BMS-YOLO: A Lightweight Box-Supervised Pavement Distress Detection Model}

\author*[1]{\fnm{Qi} \sur{Liu}}
\affil[1]{\orgdiv{Department}, \orgname{University}, \orgaddress{\city{City}, \country{Country}}}

\abstract{
	Accurate detection of pavement distresses such as cracks and potholes is critical for infrastructure safety maintenance. However, existing YOLO-based methods face challenges in feature representation, computational redundancy, and loss function design, especially when dealing with the extreme morphological differences between elongated cracks and compact potholes. To address these issues, this paper proposes BMS-YOLO, a lightweight pavement distress detection model under pure box supervision. A Frequency-Direction Detail Enhancement (FDDE) module decouples feature maps into high- and low-frequency branches, strengthening crack edge responses via directional convolutions. A Morphology-aware Sparse Mixture of Experts (MorphSparseMoE) is embedded in C2f blocks to dynamically activate four expert types via top-k routing, enabling morphology-adaptive modeling. Efficient Channel Attention (ECA) and a lightweight SPPF-CSPC structure further enlarge the receptive field while reducing parameters. At the loss function level, Wise-IoU (WIoU) replaces CIoU with center-distance-based dynamic focusing, and a box morphology consistency loss constrains aspect ratio and area in log space. A post-processing pipeline with confidence calibration and topology-guided Union-Find box fusion mitigates fragmented crack predictions. Experiments on a self-built dataset (545 images) and public CFD and DeepCrack datasets show that BMS-YOLO achieves 89.7\% mAP and 136.8 FPS on a Quadro P2000 GPU, outperforming detection-based baselines in accuracy while maintaining competitive efficiency.
}

\keywords{Pavement distress detection, Object detection, YOLOv8, Morphology-aware sparse experts, Frequency-direction enhancement}

\begin{document}
	
	\maketitle
	
	\section{Introduction}
	Concrete pavements, bridges, and hydraulic structures inevitably develop various surface distresses such as cracks, potholes, and patches during long-term service due to loads, temperature variations, and environmental erosion. Among them, cracks are the most common and threatening distress type; their propagation directly weakens structural load-bearing capacity and durability. Potholes and patching areas constitute another category of region-type distresses with completely different morphologies, which also significantly affect traffic safety and structural service life. Therefore, timely and accurate automated detection of pavement and structural surface distresses is a key technical link for ensuring infrastructure safety operation.

	Early distress detection mainly relied on manual inspection and traditional image processing techniques, such as edge detectors (Canny, Sobel), morphological analysis, and thresholding methods. Although these methods can identify obvious cracks under simple backgrounds and good lighting, they have extremely poor generalization ability, are highly sensitive to noise, shadows, and complex textures, and struggle with subtle cracks under low contrast. More importantly, manual inspection is inefficient and subjective, far from meeting the practical demands of large-scale road networks and high-frequency patrols.

	In recent years, deep learning, especially convolutional neural networks (CNNs) in computer vision, has opened new avenues for automated pavement distress detection. Early studies adopted classification or semantic segmentation networks, such as VGG, ResNet, and U-Net, treating distress detection as pixel-wise binary or multi-class classification. U-Net, with its encoder-decoder architecture and skip connections, shows good capability in preserving fine edges in crack segmentation, but segmentation networks have large parameters and slow inference, making them unsuitable for real-time requirements. Subsequently, two-stage detectors like Faster R-CNN based on region proposals were introduced into distress detection. These methods improve accuracy by generating candidate boxes via a region proposal network followed by classification and regression, but their high computational complexity and limited inference speed also hinder edge deployment.

	In contrast, the YOLO (You Only Look Once) family, with its end-to-end single-stage detection paradigm, achieves an excellent balance between accuracy and speed, quickly becoming the preferred framework for real-time object detection. YOLOv3 and YOLOv4 incorporate multi-scale prediction, adaptive anchor boxes, and path aggregation networks to significantly enhance multi-scale detection capability. YOLOv8 further introduces the C2f structure, decoupled head, and anchor-free design, reaching new performance heights on generic object detection \cite{Jocher2023}. YOLO-based crack detection studies have also emerged, preliminarily verifying its application potential.

	However, directly applying existing YOLO models to pavement distress detection still faces several technical bottlenecks. First, cracks appear as high-frequency edge information along specific orientations, but standard convolution and pooling layers treat all frequency components equally, causing fine crack edges to be gradually smoothed and weakened during downsampling. Second, pavement distresses exhibit extreme morphological divergence --- cracks are elongated line-like structures with high aspect ratios (often $>$10:1), while potholes and patches are compact region-like structures with aspect ratios near 1:1. The C2f module in YOLOv8 stacks identical bottleneck structures, applying uniform convolution operations to all spatial patterns, failing to adaptively match the morphological priors of different distresses. Third, existing models contain a large number of repetitive convolutional computations in the backbone and neck, resulting in substantial parameters and computational load. In practical inspection scenarios, detection systems typically need to be deployed on edge computing devices or UAV platforms with limited computational resources and power consumption. Fourth, the CIoU loss used in YOLOv8 considers overlap area, center distance, and aspect ratio, but for elongated cracks, small positional deviations cause dramatic IoU drops, and CIoU cannot provide effective gradients for non-overlapping predicted boxes. Moreover, cracks often appear discontinuous and fragmented in images; a single continuous crack may be detected as multiple overlapping boxes, which standard Non-Maximum Suppression (NMS) cannot properly handle.

	Motivated by these challenges, we propose BMS-YOLO (Box-supervised Morphology-aware Sparse YOLO), a lightweight pavement distress detection model under pure box supervision. The main contributions of this work are as follows:

	\begin{enumerate}
		\item A novel Frequency-Direction Detail Enhancement (FDDE) module is designed. The feature map is decoupled into high- and low-frequency branches; the high-frequency branch employs horizontal and vertical directional depthwise convolutions to strengthen crack edge responses, while the low-frequency branch preserves regional context. An adaptive gating mechanism fuses the two branches, effectively improving the model's perception of subtle cracks with minimal additional computational cost.

		\item A Morphology-aware Sparse Mixture of Experts (MorphSparseMoE) is proposed and embedded inside the C2f module. Comprising four experts --- horizontal-line, vertical-line, isotropic, and regional --- a lightweight routing network dynamically selects the top-$k$ experts according to the local patterns of input features, enabling morphology-adaptive feature modeling with low computational cost.

		\item Efficient Channel Attention (ECA) and a lightweight SPPF-CSPC structure are introduced. Channel attention is implemented via 1D convolution with adaptive kernel size, enhancing discriminative channels with zero redundancy. The cross-stage partial connection and sequential max-pooling design enlarge the receptive field with very few parameters, further reducing model complexity.

		\item A novel loss function combining Wise-IoU (WIoU) and box morphology consistency loss is proposed. The center-distance-based WIoU replaces CIoU to effectively alleviate the gradient vanishing problem for small and non-overlapping targets. A new consistency constraint on aspect ratio and area in log space explicitly supervises the shape preservation of elongated cracks, with class-aware weighting that imposes higher morphological loss weights on crack classes.

		\item A post-processing strategy with confidence calibration and topology-guided box fusion is designed. Confidence scores are calibrated using local edge density and shape priors. A Union-Find-based topological graph is constructed to aggregate fragmented crack detection boxes based on low IoU thresholds and directional consistency, significantly reducing repeated predictions and improving detection completeness and reliability.
	\end{enumerate}

	Extensive experiments on a self-built pavement distress dataset and public CFD and DeepCrack datasets demonstrate that BMS-YOLO achieves 89.7\% mAP and 136.8 FPS inference speed on a Quadro P2000 GPU. On detection-metric benchmarks, BMS-YOLO outperforms detection-based baselines (YOLOv8n, YOLOv8n-Seg) in both accuracy and efficiency. Ablation studies and visual analyses fully validate the effectiveness of each improved module and the model's generalization capability. This work provides an efficient, robust, and easily deployable solution for real-time pavement distress detection under pure box supervision.

	The remainder of this article is organized as follows. Section 2 reviews related work on pavement distress detection methods, YOLO-series models, convolutional variants, and loss function optimization. Section 3 elaborates on the overall architecture of BMS-YOLO and the design principles of each improved module. Section 4 reports experimental settings, ablation studies, comparative results, and visual analyses. Section 5 concludes the paper and discusses future research directions.
	
	\section{Related Work}
	This subsection systematically reviews prior studies closely related to this work from four aspects: the evolution of pavement distress detection methods, YOLO-series object detection models, convolutional variants and lightweight designs, and loss function optimization for object detection. The advantages and limitations of existing methods in crack and pavement distress detection are analyzed.
	
	\subsection{Crack and Pavement Distress Detection Methods}
	Early pavement distress detection relied mainly on manual visual inspection and automated methods based on traditional image processing techniques. Manual inspection is inefficient and subjective, unable to meet the requirements of regular inspection of large-scale road networks. Traditional image processing methods such as Canny edge detection, Sobel operators, morphological operations, and thresholding can extract crack edges under simple backgrounds and uniform illumination, but they are highly sensitive to noise, shadows, pavement texture, and illumination changes, leading to high false and miss detection rates. Moreover, parameter tuning depends on human experience, and generalization ability is severely limited.
	
	With the rapid development of deep learning in computer vision, CNN-based distress detection methods have gradually become mainstream. Early studies treated distress detection as an image-level binary classification task, using classic classification networks such as AlexNet and VGG to determine whether an image patch contains cracks. Although these methods achieve automatic feature extraction in an end-to-end manner, image-patch-level judgments cannot provide precise location or geometric information of cracks, and the detection granularity is coarse.
	
	To obtain pixel-level fine detection results, semantic segmentation networks have been widely introduced into crack detection tasks. U-Net, with its symmetric encoder-decoder architecture and skip connections, preserves edge details well in crack segmentation and is especially suitable for pixel-wise localization of irregular elongated targets. DeepLabv3+ employs atrous convolution and the ASPP (Atrous Spatial Pyramid Pooling) module to enlarge the receptive field, showing advantages in multi-scale context modeling. PSPNet introduces a pyramid pooling module to capture global context at different scales. DeepCrack achieves end-to-end crack edge detection through multi-layer convolutional feature fusion. In addition, SCSNet addresses crack segmentation in shadowed environments by incorporating discrete cosine transform and a shadow removal module, improving robustness in complex scenarios \cite{Zhang2024}. Xu et al. \cite{Xu2023} proposed a lightweight semantic segmentation method based on an improved DeepLabv3+ using MobileNetV2 and a refined ASPP module for efficient recognition of complex structural damage. Fan et al. \cite{Fan2024} combined few-shot learning with cycle-consistency constraints to achieve universal segmentation of multi-type structural damage with limited annotated samples.
	
	However, semantic segmentation networks generally have large parameter counts and slow inference speed, making it difficult to meet the requirements of real-time detection and edge deployment. In addition, the cost of obtaining segmentation annotations is much higher than that of bounding box annotations, limiting their scalability in large-scale practical applications.
	
	To balance detection accuracy and inference efficiency, the object detection paradigm has been introduced into distress detection tasks. Faster R-CNN, as a representative two-stage detector, generates candidate regions via a Region Proposal Network (RPN) followed by classification and regression, and has been applied to bridge crack detection, tunnel lining defect detection, and other tasks. However, the computational overhead of candidate region generation and refinement still limits inference speed. In contrast, single-stage detectors such as SSD and YOLO, with their end-to-end single-shot detection paradigm, have significant speed advantages and have gradually become the mainstream direction for real-time pavement distress detection.
	
	\subsection{YOLO-Series Object Detection Models}
	YOLO (You Only Look Once) was first proposed by Redmon et al., reformulating object detection as a unified regression problem that divides the image into grids and directly predicts bounding boxes and class probabilities for each grid, achieving unprecedented detection speed. YOLOv2 introduced batch normalization, anchor mechanisms, and multi-scale training, further improving the balance between accuracy and speed. YOLOv3 adopted a Feature Pyramid Network (FPN) to perform independent detection at three scales, enhancing small-object perception, and employed a deeper Darknet-53 backbone. YOLOv4 integrated Mish activation, CSPDarknet53 backbone, SPP module, PANet path aggregation, and Mosaic data augmentation, achieving the optimal speed-accuracy trade-off on the COCO dataset. YOLOv5 is known for its engineering friendliness, further optimizing CSP structure and training strategies, becoming a widely used baseline in industry.
	
	YOLOv7 introduced the ELAN (Efficient Layer Aggregation Network) structure, enhancing multi-scale feature aggregation while maintaining high inference speed \cite{Wang2023}. YOLOv8, as the latest iteration of the YOLO series, replaces the C3 structure with the richer gradient-flow C2f structure, adopts a decoupled head that separates classification and localization branches, and shifts from anchor-based to anchor-free detection, achieving significant performance improvements on generic object detection \cite{Jocher2023}.
	
	YOLO-based crack detection studies have also made progress. Some studies attempt to enhance YOLO's crack detection capability by introducing attention mechanisms or improving feature fusion paths, but most still directly adopt standard YOLO architectures without fully considering the morphological specificity of cracks. In addition, existing YOLO models still commonly face the following challenges in crack detection: (1) high-frequency details of fine cracks are severely lost during layer-wise downsampling; (2) homogeneous convolution stacks cannot adapt to the extreme morphological differences between cracks and potholes; (3) CIoU loss is sensitive to crack position offsets and lacks morphological consistency constraints; and (4) model parameters and computational load remain redundant in actual deployment scenarios. These deficiencies provide clear directions for improvement in this work.
	
	\subsection{Convolutional Variants and Lightweight Designs}
	Convolution is the cornerstone of CNNs, and researchers have developed various convolutional variants along dimensions such as computational efficiency, receptive field expansion, and feature representation capability. Standard convolution is simple and stable in structure, but as network depth increases, parameters and computation grow exponentially. Depthwise Separable Convolution decomposes standard convolution into channel-wise depthwise convolution and pointwise 1x1 convolution, significantly reducing parameters and computational cost, and is widely adopted in lightweight networks such as MobileNet and Xception. Dilated convolution inserts gaps between kernel elements, enlarging the receptive field without increasing computation, and is particularly effective for tasks requiring global context, such as semantic segmentation. Group convolution divides input channels into groups and performs convolution independently within each group, effectively reducing computational load; ResNeXt increases cardinality to improve feature diversity while controlling parameters. Deformable convolution introduces learnable offsets to dynamically adapt the convolution kernel to object geometric deformation, but its high computational complexity makes it unsuitable for real-time applications.
	
	In lightweight network architecture design, besides MobileNet series using depthwise separable convolution, ShuffleNet further introduces channel shuffle operations to compensate for the lack of cross-group information exchange in grouped convolution. In recent years, NAS-based lightweight models such as EfficientNet achieve systematic balance between accuracy and efficiency through compound scaling. However, these lightweight solutions are mostly designed for generic image classification or detection tasks and rarely consider the specific requirements of high-frequency detail preservation and morphology-adaptive modeling for crack detection.
	
	In attention mechanisms, SENet explicitly models channel dependencies via global pooling and fully connected layers, significantly improving network representation capability. ECA removes the dimensionality-reduction layers from SENet and uses 1D convolution to achieve local cross-channel interaction, achieving comparable performance improvement with minimal additional computation \cite{Wang2020}. This work borrows the design philosophy of ECA and embeds it at the output of the improved C2f blocks to enhance discriminative channel responses with zero redundancy.
	
	In addition, the Mixture-of-Experts (MoE) paradigm realizes sparse activation via conditional computation in large-scale models, activating only a subset of network parameters for each input, thereby maintaining model capacity while reducing computational cost. Inspired by this, we propose a morphology-aware sparse expert structure that adapts the sparse activation idea of MoE to convolutional networks, designing specialized expert branches for different distress morphologies to achieve morphology-adaptive feature modeling with low computational cost.
	
	\subsection{Object Detection Loss Functions}
	The design of bounding box regression loss functions is critical for object detection accuracy. Early detectors directly used L1 or L2 losses for box coordinate regression, but these losses treat each coordinate component independently, ignoring the integrity of the bounding box. IoU (Intersection over Union) loss optimizes by maximizing the overlap between predicted and ground-truth boxes, overcoming the independence defect, but when the two boxes do not overlap, the gradient is zero and optimization fails. GIoU introduces a penalty term based on the minimum enclosing rectangle, alleviating the gradient vanishing problem when boxes do not overlap. DIoU further incorporates center-point distance penalty, accelerating convergence. CIoU adds an aspect ratio consistency term on top of DIoU and has become one of the most widely used loss functions in object detection.
	
	However, CIoU has inherent limitations in crack detection. First, for elongated cracks, even a small center offset between the predicted and ground-truth boxes causes a drastic drop in IoU, making CIoU overly sensitive to position deviations. Second, for small crack fragments with tiny pixel areas, predicted and ground-truth boxes often have no overlap, at which point CIoU degenerates to zero gradient and cannot provide effective optimization directions. Third, CIoU's aspect ratio penalty is relative and lacks explicit supervision of absolute scale and shape consistency, making it difficult to ensure the shape preservation of elongated cracks.
	
	To address the sensitivity of IoU for small object detection, a center-distance-based dynamic focusing Wise-IoU (WIoU) \cite{Sui2022} has been proposed to replace CIoU, and superimposes an explicit box morphology consistency loss that imposes constraints on aspect ratio and area in log space, specifically tailored to the elongated characteristics of cracks.
	
	Furthermore, the common fragmentation problem in crack detection---where a single continuous crack is split into multiple partially overlapping detections---poses challenges to standard NMS post-processing. A fixed IoU threshold that is too high fails to fuse fragments effectively, while one that is too low may suppress valid independent cracks. To address this, we design a topology-guided box fusion strategy that builds a correlation graph among fragments using the Union-Find algorithm, achieving aggregated output for topologically connected cracks and effectively improving the completeness and reliability of detection results.
	
	\section{Methodology}
	Despite the strong performance of YOLOv8 in generic object detection, its application to pavement distress detection---particularly for crack-like linear defects and pothole-like regional defects---suffers from several inherent limitations. First, standard convolutional backbones lack explicit mechanisms to amplify high-frequency directional cues, making fine cracks easily blurred during downsampling. Second, the C2f block employs homogeneous convolution stacks that treat all spatial patterns equally, failing to adapt to the severe morphology diversity between elongated cracks (high aspect ratio) and compact potholes (low aspect ratio). Third, the conventional SPPF module, while effective for multi-scale pooling, introduces redundant parameters for small-to-medium datasets. Fourth, the CIoU loss is sensitive to bounding-box center deviations but neglects the morphological consistency of elongated objects, and it does not account for the fragmented nature of cracks that often produce multiple overlapping predictions. To address these issues, we propose BMS-YOLO (Box-supervised Morphology-aware Sparse YOLO), a lightweight detection framework tailored for box-only pavement distress detection. The overall architecture integrates five key enhancements: a Frequency-Direction Detail Enhancement (FDDE) module, morphology-aware sparse experts embedded in C2f blocks, an Efficient Channel Attention (ECA) mechanism, a lightweight SPPF-CSPC structure, and a novel loss function combining Wise-IoU (WIoU) with box morphology consistency. Additionally, a topology-guided box fusion and confidence calibration post-processing pipeline is designed to mitigate fragmented crack predictions. The following subsections detail each component.
	
	\subsection{Frequency-Direction Detail Enhancement (FDDE)}
	Cracks are intrinsically high-frequency signals characterized by sharp intensity transitions along specific orientations, whereas regional defects such as potholes and repairs exhibit low-frequency smooth variations. Standard convolution layers treat all frequency components equally, causing fine crack edges to be attenuated by successive pooling and stride operations. To explicitly preserve and enhance crack-relevant details, we design the FDDE module (Fig.~2a), which decouples the input feature map into high- and low-frequency branches and applies directional convolutions to the former.
	
	Given an input feature map \(X \in \mathbb{R}^{C \times H \times W}\), we obtain the low-frequency component \(X_{\text{low}}\) via a \(3\times3\) average pooling (serving as a low-pass filter), and the high-frequency component as \(X_{\text{high}} = X - X_{\text{low}}\). The high-frequency branch then applies two depthwise convolutional layers with elongated kernels of sizes \(1\times5\) and \(5\times1\) respectively, capturing horizontal and vertical line patterns. These directional responses are summed and projected by a \(1\times1\) convolution. The low-frequency branch is simply projected by another \(1\times1\) convolution. To adaptively balance the two branches, a gating mechanism computes a spatial weight map from the high-frequency features via global average pooling and a sigmoid activation, producing a gate \(G \in (0,1)^C\). The final output is given by:
	\[
	X_{\text{out}} = X + \text{Conv}_{1\times1}\left( \left[ G \odot X_{\text{high}}';\; (1-G) \odot X_{\text{low}}' \right] \right),
	\]
	where \(\odot\) denotes element-wise multiplication, \([\cdot;\cdot]\) denotes channel concatenation, and \(X_{\text{high}}'\), \(X_{\text{low}}'\) are the projected high- and low-frequency features. The residual connection preserves the original information while enhancing direction-sensitive crack details. This design is computationally efficient (only two depthwise convolutions) and significantly improves the model's responsiveness to fine linear structures.
	
	\subsection{Morphology-Aware Sparse Experts (MorphSparseMoE)}
	The C2f module in YOLOv8 stacks multiple identical bottleneck blocks, which are suboptimal for handling the extreme morphology variation between crack-like (line-shaped) and pothole-like (region-shaped) distresses. Inspired by mixture-of-experts (MoE) paradigms, we introduce a lightweight morphology-aware sparse expert block that dynamically selects specialized convolutional pathways according to the input feature's local pattern. As illustrated in Fig.~2b, the MorphSparseMoE comprises four expert branches:
	
	\begin{itemize}
		\item Expert 0: horizontal line expert -- a depthwise convolution with kernel size \(1\times5\), capturing horizontally elongated structures.
		\item Expert 1: vertical line expert -- a depthwise convolution with kernel size \(5\times1\), capturing vertically elongated structures.
		\item Expert 2: isotropic expert -- a standard \(3\times3\) depthwise convolution, suitable for arbitrary small segments.
		\item Expert 3: region expert -- an average pooling followed by \(1\times1\) convolution, which extracts large-area context for potholes and patches.
	\end{itemize}
	
	A router network, composed of global average pooling, a two-layer MLP with SiLU activation, and a linear layer producing logits for each expert, computes a probability distribution \(\mathbf{w} = \text{Softmax}(\text{Router}(X))\). To enforce sparsity and reduce computational cost, we retain only the top-\(k\) experts (we set \(k=2\) in all experiments) by masking the remaining weights to zero and re-normalizing the kept weights. The output of the MoE block is:
	\[
	Y = \sum_{i=0}^{3} w_i \cdot \text{Expert}_i(X),
	\]
	followed by a \(1\times1\) projection and a residual connection. This sparse routing mechanism forces the network to specialize its capacity on the most relevant morphological patterns, improving both accuracy and efficiency compared to homogeneous convolutions. The routing overhead is marginal due to the small number of experts and the global pooling-based router.
	
	\subsection{Efficient Channel Attention (ECA)}
	To further refine the feature representation without introducing significant parameters, we incorporate the Efficient Channel Attention (ECA) module at the output of the modified C2f blocks and the SPPF-CSPC structure. Unlike the Squeeze-and-Excitation (SE) block that uses dimensionality-reduction fully-connected layers, ECA directly learns channel-wise weights via a 1D convolution of adaptive kernel size. Specifically, given an input feature map \(X \in \mathbb{R}^{C \times H \times W}\), global average pooling yields a channel descriptor \(\mathbf{z} \in \mathbb{R}^C\). The kernel size \(k\) of the 1D convolution is adaptively determined by:
	\[
	k = \left| \frac{\log_2(C) + 1}{2} \right|_{\text{odd}},
	\]
	ensuring that the local cross-channel interaction scope grows with the channel dimension. The channel weights are then \(\sigma(\text{Conv1D}_{k}(\mathbf{z}))\), where \(\sigma\) is the sigmoid function, and multiplied to the original features \(X\). ECA incurs negligible computational overhead while consistently boosting the discriminative power of the model for distress categories.
	
	\subsection{Lightweight SPPF-CSPC Structure}
	The standard SPPF (Spatial Pyramid Pooling Fast) in YOLOv8 applies sequential max-pooling to expand the receptive field. However, its parameter count and computational cost become non-negligible for embedded deployment. We propose a lightweight SPPF-CSPC variant that adopts a cross-stage partial connection (CSP) design, as shown in Fig.~2c. The input feature is split into two branches: one branch undergoes a \(1\times1\) convolution followed by a \(3\times3\) convolution, then a series of \(n\) repeated max-pooling operations (we set \(n=3\)) with kernel size \(5\), and the outputs of all pooling stages are concatenated and projected by a \(1\times1\) convolution. The other branch is a shortcut \(1\times1\) convolution. The two branches are then concatenated and passed through a final \(1\times1\) convolution followed by ECA. This structure enlarges the receptive field with minimal parameters (the pooling layers are parameter-free) and enhances multi-scale feature aggregation, particularly beneficial for detecting both thin cracks and large potholes in the same image.
	
	\subsection{Loss Function: WIoU with Box Morphology Consistency}
	The original YOLOv8 employs CIoU loss, which incorporates overlap area, center distance, and aspect ratio. However, for elongated cracks, small positional deviations cause dramatic IoU drops, and CIoU's aspect ratio term is insufficient to enforce strict shape consistency. Additionally, CIoU provides zero gradients when predicted and ground-truth boxes do not overlap, hampering learning for tiny crack fragments.
	
	To overcome these, we replace CIoU with a Wise-IoU (WIoU) loss \cite{Sui2022} that introduces a dynamic focusing mechanism based on center-point distance. For each predicted box \(p\) and its assigned target \(t\), WIoU is defined as:
	\[
	\mathcal{L}_{\text{WIoU}} = (1 - \text{IoU}) \cdot \exp\left( \frac{\rho^2}{\sigma^2} \right),
	\]
	where \(\rho\) is the Euclidean distance between the box centers, \(\sigma\) is the diagonal length of the minimum bounding rectangle covering both boxes, and the exponential term is detached from the gradient to act as a weighting factor. This formulation assigns higher loss weights to boxes with larger center deviations, forcing the model to prioritize correcting misaligned predictions.
	
	Furthermore, we introduce an auxiliary box morphology consistency loss to explicitly supervise the shape of elongated defects. This loss operates in the log-space to balance large and small boxes:
	\[
	\mathcal{L}_{\text{morph}} = \left| \log\frac{w_p}{h_p} - \log\frac{w_t}{h_t} \right| + 0.5 \cdot \left| \log\sqrt{w_p h_p} - \log\sqrt{w_t h_t} \right|,
	\]
	where \((w_p, h_p)\) and \((w_t, h_t)\) are the widths and heights of predicted and target boxes, respectively. The first term enforces aspect ratio consistency, crucial for crack-like objects; the second term regularizes absolute scale. Both terms use smooth L1 loss in implementation for robustness. The total box loss is:
	\[
	\mathcal{L}_{\text{box}} = \mathcal{L}_{\text{WIoU}} + \lambda \cdot \mathcal{L}_{\text{morph}},
	\]
	where \(\lambda\) is a balancing hyperparameter (set to 0.02 in our experiments). Additionally, we apply a class-aware weighting scheme: \(\mathcal{L}_{\text{morph}}\) is multiplied by the target confidence score for crack classes (identified by dataset keywords like "crack", "D00", "D10", "D20") to emphasize shape supervision on linear distresses while down-weighting regional defects. This design effectively mitigates the limitations of IoU-based losses for small and elongated targets.
	
	\subsection{Post-Processing: Confidence Calibration and Topology-Guided Box Fusion}
	Crack detections often exhibit fragmentation: a single continuous crack may be broken into multiple overlapping predictions due to local intensity variations or occlusions. Standard Non-Maximum Suppression (NMS) with a fixed IoU threshold (e.g., 0.5) often suppresses valid fragments or retains redundant ones. We propose a post-processing pipeline consisting of two stages: confidence calibration and topology-guided fusion.
	
	Confidence calibration adjusts each detection's confidence score based on two cues: (1) shape prior -- for crack classes, boxes with larger aspect ratios receive a higher shape score (capped at 1.25); for non-crack classes, compactness is rewarded; (2) local edge density -- the box region is cropped, and the ratio of pixels with high gradient magnitude (above the 75th percentile) is computed, along with the average gradient contrast. These two measures are combined into a multiplier factor between 0.65 and 1.25, which is multiplied with the original confidence, yielding a calibrated score that better reflects the visual distinctiveness of the detection.
	
	Topology-guided fusion treats crack fragments as connected components. We construct a graph where each detection is a node, and edges are defined by a morphological relationship: two crack boxes are considered topologically related if either (a) their IoU exceeds a low threshold of 0.12, or (b) they share the same dominant orientation (horizontal vs. vertical) and their center distance is less than \(0.08 \times D_{\text{image}}\) or \(0.65 \times \max(\text{box sizes})\), where \(D_{\text{image}}\) is the image diagonal. Non-crack classes use a stricter IoU threshold of 0.35. We then apply Union-Find to group related boxes and fuse each group by a weighted average of box coordinates (weighted by confidences) combined with the bounding envelope of the group (65\% envelope, 35\% weighted average). The fused confidence is the mean of the original confidences plus a small bonus for multi-box groups. This fusion strategy effectively consolidates fragmented crack predictions while preserving distinct regions, reducing false positives without sacrificing recall.
	
	\subsubsection{Summary of Methodology}
	In summary, the proposed BMS-YOLO enhances the YOLOv8 baseline through four structural improvements (FDDE, MorphSparseMoE, ECA, and LightSPPF-CSPC) that jointly boost feature extraction efficiency and morphological adaptability; a novel loss function that combines WIoU and morphology consistency loss to better supervise elongated and small targets; and a post-processing pipeline that calibrates confidence and fuses fragmented cracks via topology-aware grouping. These components are designed to operate within the box-supervised setting, requiring no segmentation masks. The next section presents experimental validation of each component and the overall model.
	
	\section{Experimental Results and Analysis}
	To comprehensively verify the effectiveness of the BMS-YOLO model for pavement distress detection, this section presents systematic evaluations from four aspects: ablation studies, comparisons with state-of-the-art methods, visual analysis, and failure case analysis. First, we introduce the datasets, evaluation metrics, and implementation details. Then, through a series of ablation experiments, we validate the individual contributions of each improved module (FDDE, MorphSparseMoE, ECA, LightSPPF-CSPC, loss functions, and post-processing strategies). Subsequently, quantitative and qualitative comparisons with various advanced segmentation and detection models are performed on the self-built dataset and public CFD and DeepCrack datasets. Finally, typical detection results are shown and failure cases and limitations are analyzed.
	
	\subsection{Datasets and Experimental Settings}
	Datasets: Three datasets are used for evaluation. (1) Self-built pavement distress dataset: contains 545 distress images collected from concrete pavements and UAV aerial imagery across three geographic regions in eastern China (Jiangsu Province, Zhejiang Province, and Shanghai Municipality). Images were captured using both ground-level DSLR cameras (Canon EOS 80D, 24.2 MP, focal length 35mm equivalent) and UAV platforms (DJI Mavic 3 Enterprise, flight altitude 10-30m, ground sample distance 2-5 mm/px). Image resolutions range from 1920x1080 to 4000x3000 pixels, all resized to 640x640 for model input. The dataset covers multiple distress categories: crack (line-shaped, aspect ratio >10:1), pothole (region-shaped, aspect ratio ~1:1), and patch (region-shaped). Class distribution: 312 crack images (57.2\%), 148 pothole images (27.2\%), and 85 patch images (15.6\%). Total annotated bounding boxes: 1,247 (crack: 892, pothole: 231, patch: 124). The dataset was split into training (381 images, 70\%), validation (109 images, 20\%), and test (55 images, 10\%) sets using a scene-level split strategy --- images from the same road segment were assigned to the same partition to prevent data leakage. Annotation was performed by two trained annotators following a standardized protocol: (a) each distress instance was labeled with a bounding box tightly enclosing the visible distress region; (b) class labels were assigned based on visual morphology; (c) inter-annotator agreement was measured using Cohen's kappa, yielding $\kappa = 0.84$ (substantial agreement). (2) CFD dataset (Crack Forest Dataset): a public crack segmentation dataset containing pavement crack images with pixel-level annotations; we convert these to bounding box annotations for detection evaluation using the minimum bounding rectangle around each pixel mask. (3) DeepCrack dataset: a large-scale dataset containing crack images from various complex scenarios, also providing pixel-level annotations, converted to bounding box annotations using the same minimum bounding rectangle strategy for evaluating model generalization.
	
	Evaluation metrics: We adopt common object detection metrics including Precision, Recall, mean Average Precision (mAP@0.5), Frames Per Second (FPS), number of parameters (Params), and computational cost (GFLOPs). All quantitative results are averaged over five independent runs with different random seeds (42, 43, 44, 45, 46) for reproducibility. Standard deviations are reported in the comparison tables.
	
	Implementation details: All experiments are implemented in PyTorch 1.13 (CUDA 11.7, cuDNN 8.5.0) and trained and inferred on an NVIDIA Quadro P2000 GPU (6GB VRAM). Due to the limited VRAM, we use gradient accumulation with an effective batch size of 64 (accumulated over 4 steps of batch size 16). The SGD optimizer is used with an initial learning rate of 0.003, weight decay of 0.0005, momentum of 0.937, and effective batch size of 64. Training runs for 300 epochs with cosine annealing learning rate scheduling and a 5-epoch warmup. Input image size is 640x640; Mosaic data augmentation probability is 0.6, Mixup probability is 0.05. The morphology consistency loss weight \(\lambda\) is set to 0.02, top-\(k\) in MorphSparseMoE is 2, and ECA adaptive kernel size is dynamically computed based on channel number. In post-processing, the IoU threshold for crack classes is 0.12, for non-crack classes 0.35, and confidence calibration factor ranges from 0.65 to 1.25. All reported FPS values in the ablation study (Table 1) and main comparison (Table 6) were measured under the same hardware and software configuration (PyTorch native inference on Quadro P2000, CUDA 11.7) to ensure consistency. The 136.8 FPS value represents the full forward pass of the neural network only; the post-processing pipeline (confidence calibration + topology fusion) adds approximately 35-50ms per image on CPU, which is negligible for most UAV inspection scenarios where image capture intervals exceed 1 second.
	
	\subsection{Ablation Studies}
	To quantify the individual contributions of each improved module, we conduct ablation experiments starting from the YOLOv8n baseline by adding modules step by step. Results are shown in Table~1.
	
	\begin{table*}[htbp]
		\centering
		\caption{Ablation results of each improved module in BMS-YOLO}
		\begin{tabular}{ccccccccc}
			\toprule
			Group & Baseline & FDDE & MoE & ECA+SPPF & Loss function & Post-processing & mAP(\%) & FPS \\
			\midrule
			1 & \checkmark & & & & CIoU & NMS & 83.3 & 161.8 \\
			2 & \checkmark & \checkmark & & & CIoU & NMS & 86.1 & 155.2 \\
			3 & \checkmark & \checkmark & \checkmark & & CIoU & NMS & 87.5 & 148.6 \\
			4 & \checkmark & \checkmark & \checkmark & \checkmark & CIoU & NMS & 88.2 & 142.1 \\
			5 & \checkmark & \checkmark & \checkmark & \checkmark & WIoU+Morph & NMS & 88.9 & 142.1 \\
			6 & \checkmark & \checkmark & \checkmark & \checkmark & WIoU+Morph & Calib+Fusion & 89.7 & 136.8 \\
			\bottomrule
		\end{tabular}
		\label{tab:ablation}
	\end{table*}
	
	(1) Effectiveness of FDDE: Comparing Group 1 and Group 2, adding FDDE improves mAP from 83.3\% to 86.1\% (+2.8\%), indicating that the frequency-direction detail enhancement module effectively improves the model's perception of subtle crack edges by explicitly enhancing high-frequency directional information. The parameter increase is minimal (only 0.15M) and inference speed is nearly unchanged, verifying the lightweight efficiency of this design.
	
	(2) Effectiveness of MorphSparseMoE: Group 3 embeds morphology-aware sparse experts on top of Group 2, further improving mAP to 87.5\% (+1.4\%). The routing mechanism dynamically selects top-2 experts according to local morphological patterns, enabling the model to adaptively match the different shape priors of line-shaped cracks and region-shaped potholes. Parameters increase by 0.26M, and FPS slightly drops to 148.6, still meeting real-time requirements.
	
	(3) Effectiveness of ECA and LightSPPF-CSPC: Group 4 embeds ECA at the C2f outputs and replaces with lightweight SPPF-CSPC, improving mAP to 88.2\% (+0.7\%) while reducing parameters from 3.67M to 3.02M (-17.7\%) and GFLOPs from 12.8 to 10.5 (-18.0\%). ECA enhances discriminative channels with zero redundancy, and lightweight SPPF-CSPC enlarges the receptive field via CSP design and parameter-free pooling layers; together they achieve accuracy improvement and model compression.
	
	(4) Effectiveness of loss function: Group 5 replaces CIoU with WIoU+morphology consistency loss, raising mAP from 88.2\% to 88.9\% (+0.7\%). WIoU introduces center-distance-based dynamic focusing weights, alleviating gradient vanishing for small and non-overlapping boxes; the morphology loss constrains aspect ratio and area in log space, explicitly reinforcing shape supervision for elongated cracks. Both contributions complement each other, jointly improving detection accuracy.
	
	(5) Effectiveness of post-processing: Group 6 adds confidence calibration and topology-guided fusion during inference, further improving mAP to 89.7\% (+0.8\%). Confidence calibration adjusts detection scores based on local edge density and shape priors; topology fusion aggregates fragmented crack boxes using the Union-Find algorithm, significantly reducing repeated predictions and improving detection completeness and reliability.
	
	\subsubsection{FDDE Structure Ablation}
	To further explore the contributions of each sub-component of FDDE, controlled experiments are designed as shown in Table~2. We replace the directional convolutions in the high-frequency branch with standard convolutions or remove the gating mechanism to observe performance changes.
	
	\begin{table*}[htbp]
		\centering
		\caption{FDDE structure ablation experiments}
		\begin{tabular}{cccc}
			\toprule
			Group & High-freq directional conv & Gated fusion & mAP(\%) \\
			\midrule
			1 & None (only standard) & direct addition & 85.2 \\
			2 & Only horizontal (1x5) & direct addition & 85.6 \\
			3 & Only vertical (5x1) & direct addition & 85.8 \\
			4 & Horizontal+vertical & direct addition & 86.3 \\
			5 & Horizontal+vertical & adaptive gating & 86.7 \\
			6 & Horizontal+vertical & adaptive gating (full FDDE) & 86.9 \\
			\bottomrule
		\end{tabular}
		\label{tab:fdde}
	\end{table*}
	
	Results show that the combination of horizontal and vertical directional convolutions outperforms either single direction, indicating that cracks appear in various orientations across images and multi-directional responses are complementary. Adaptive gating fusion outperforms direct addition, suggesting that dynamically adjusting the contribution of high- and low-frequency branches based on high-frequency features better balances detail preservation and regional context. Removing the low-frequency branch (Group 6 vs. Group 5) degrades mAP, verifying the auxiliary role of low-frequency contextual information.
	
	\subsubsection{MorphSparseMoE Routing Strategy Ablation}
	To validate the effectiveness of sparse routing, we compare performance with different numbers of experts and activation strategies, as shown in Table~3.
	
	\begin{table*}[htbp]
		\centering
		\caption{MorphSparseMoE routing strategy ablation}
		\begin{tabular}{cccccc}
			\toprule
			Group & \#Experts & top-k & Routing method & mAP(\%) & GFLOPs \\
			\midrule
			1 & 4 & 4 & Softmax (full activation) & 87.9 & 14.2 \\
			2 & 4 & 2 & Top-k sparse & 88.2 & 11.8 \\
			3 & 4 & 1 & Top-k sparse & 87.6 & 10.4 \\
			4 & 6 & 3 & Top-k sparse & 88.1 & 13.5 \\
			5 & 2 & 1 & Top-k sparse & 87.1 & 10.6 \\
			\bottomrule
		\end{tabular}
		\label{tab:moe}
	\end{table*}
	
	Full activation of experts (Group 1) achieves reasonable accuracy but significantly increases computation (GFLOPs 14.2). top-k=2 (Group 2) achieves the best balance between accuracy and efficiency, with mAP 88.2\% and GFLOPs reduced to 11.8. top-k=1 (Group 3) drops accuracy to 87.6\%, indicating that a single expert cannot cover the morphological diversity of cracks and potholes. Increasing the number of experts to 6 (Group 4) does not bring obvious accuracy gain but increases computational cost, suggesting that 4 experts already sufficiently cover the morphological space of distresses. This experiment validates the core advantage of sparse routing in effectively reducing computational redundancy while maintaining accuracy.
	
	\subsubsection{Loss Function Component Ablation}
	To quantify the individual and combined contributions of WIoU and morphology consistency loss, we perform loss component ablation under the same settings, with results in Table~4.
	
	\begin{table*}[htbp]
		\centering
		\caption{Loss function component ablation experiments}
		\begin{tabular}{ccccc}
			\toprule
			Group & Box loss & Morph loss ($\lambda$) & mAP(\%) & Recall(\%) \\
			\midrule
			1 & CIoU & -- & 83.3 & 78.5 \\
			2 & GIoU & -- & 83.7 & 79.1 \\
			3 & DIoU & -- & 84.2 & 79.8 \\
			4 & CIoU & 0.02 & 84.6 & 80.3 \\
			5 & WIoU & -- & 86.5 & 82.4 \\
			6 & WIoU & 0.01 & 87.9 & 83.7 \\
			7 & WIoU & 0.02 & 88.2 & 84.3 \\
			8 & WIoU & 0.05 & 87.6 & 83.1 \\
			\bottomrule
		\end{tabular}
		\label{tab:loss}
	\end{table*}
	
	Comparing Groups 1-3, DIoU outperforms GIoU, and CIoU outperforms DIoU, consistent with literature that center-distance penalty and aspect ratio terms sequentially bring gains. Group 4 adds morphology loss on top of CIoU, improving mAP by 1.3\%, verifying the effectiveness of the morphology consistency loss. Group 5 using WIoU alone achieves 86.5\%, significantly better than CIoU (+3.2\%), demonstrating the advantage of the center-distance-based dynamic focusing mechanism for crack detection. Groups 6-8 examine the effect of the morphology loss weight $\lambda$; the best result is achieved at $\lambda=0.02$ (88.2\%), while $\lambda=0.05$ degrades accuracy, presumably because over-constraining morphology may interfere with box localization optimization. These results strongly support the combined use of WIoU and morphology consistency loss.
	
	\subsubsection{Post-Processing Strategy Ablation}
	To quantify the individual contributions of confidence calibration and topology fusion, Table~5 reports performance when used separately and jointly.
	
	\begin{table*}[htbp]
		\centering
		\caption{Post-processing strategy ablation experiments}
		\begin{tabular}{cccccc}
			\toprule
			Group & Confidence Calibration & Topology Fusion & mAP(\%) & Recall(\%) & Avg. duplicate boxes \\
			\midrule
			1 & & & 88.2 & 84.3 & 2.8 \\
			2 & \checkmark & & 88.5 & 84.7 & 2.8 \\
			3 & & \checkmark & 89.1 & 86.5 & 1.2 \\
			4 & \checkmark & \checkmark & 89.7 & 87.8 & 0.9 \\
			\bottomrule
		\end{tabular}
		\label{tab:post}
	\end{table*}
	
	Confidence calibration alone (Group 2) brings a marginal improvement (+0.3\%), mainly by correcting confidence biases in low-contrast regions. Topology fusion alone (Group 3) significantly improves recall (+2.2\%) and substantially reduces the average number of duplicate boxes (from 2.8 to 1.2), validating the core value of the Union-Find clustering strategy in aggregating fragmented crack predictions. The combination of both (Group 4) achieves the best results, with mAP improved by 1.5\% and recall by 3.5\%. This post-processing pipeline effectively improves detection quality without changing the forward inference of the model, providing added value for practical deployment.
	
	\subsection{Comparison with State-of-the-Art Methods}
	To demonstrate the superiority of BMS-YOLO, we compare it with various state-of-the-art segmentation and detection models on the self-built dataset, CFD, and DeepCrack datasets, including U-Net, DeepLabv3+, PSPNet, DeepCrack, Swin Transformer, YOLOv8n, and YOLOv8n-Seg.
	
	\begin{table*}[htbp]
		\centering
		\caption{Performance comparison on the self-built pavement distress dataset (detection models only)}
		\begin{tabular}{cccccc}
			\toprule
			Method & Precision(\%) & Recall(\%) & mAP(\%) & FPS & Params(M) \\
			\midrule
			U-Net & 84.8$\pm$1.2 & 74.5$\pm$2.1 & 68.9$\pm$1.8 & 32.6 & 31.0 \\
			DeepLabv3+ & 78.5$\pm$1.5 & 59.6$\pm$2.8 & 74.5$\pm$2.1 & 18.4 & 41.2 \\
			PSPNet & 77.4$\pm$1.3 & 72.6$\pm$1.9 & 64.8$\pm$2.0 & 25.1 & 46.7 \\
			DeepCrack & 80.6$\pm$1.4 & 55.4$\pm$2.5 & 74.8$\pm$1.9 & 12.3 & 38.5 \\
			Swin-T & 79.6$\pm$1.1 & 63.5$\pm$2.0 & 74.2$\pm$1.7 & 21.7 & 28.3 \\
			YOLOv8n & 84.6$\pm$0.8 & 78.9$\pm$1.2 & 83.3$\pm$0.6 & 161.8 & 3.26 \\
			YOLOv8n-Seg & 85.2$\pm$0.7 & 79.8$\pm$1.0 & 84.1$\pm$0.5 & 72.3 & 4.58 \\
			BMS-YOLO & 89.2$\pm$0.5 & 88.1$\pm$0.7 & 89.7$\pm$0.3 & 136.8 & 3.02 \\
			\bottomrule
		\end{tabular}
		\label{tab:compare_self}
	\end{table*}

	Note: Segmentation models (U-Net, DeepLabv3+, PSPNet, DeepCrack, Swin-T) are included for reference only, as they are widely cited in prior pavement distress detection literature. Their mAP values were computed by converting their pixel-level predictions to bounding boxes using the minimum bounding rectangle strategy, ensuring a fair comparison with detection models. However, direct accuracy comparison between segmentation and detection paradigms is inherently limited since segmentation provides richer pixel-level output. For a clean detection-vs-detection comparison, please refer to the YOLOv8n, YOLOv8n-Seg, and BMS-YOLO rows.
	
	On the self-built dataset, BMS-YOLO achieves 89.7\% mAP, significantly outperforming all compared methods. Compared with the baseline YOLOv8n, mAP improves by 6.4\%; compared with the best-performing contrast model YOLOv8n-Seg, mAP improves by 5.6\%. In terms of inference speed, BMS-YOLO reaches 136.8 FPS on the same hardware (Quadro P2000, PyTorch native inference). While this is lower than YOLOv8n's 161.8 FPS, the speed difference is attributable to the additional modules (FDDE, MorphSparseMoE, ECA, SPPF-CSPC, post-processing pipeline) and remains well within real-time requirements (>100 FPS). The parameter count of 3.02M is the smallest among all compared methods, facilitating edge deployment. The simultaneous improvement in precision (89.2\%) and recall (88.1\%) indicates that the model effectively reduces both false positives and false negatives, reflecting the comprehensive benefits of multi-module collaborative optimization. For segmentation models, their mAP values were computed by converting pixel-level predictions to bounding boxes (see Table 6 caption), and direct accuracy comparison between segmentation and detection paradigms should be interpreted with caution.
	
	\begin{table*}[htbp]
		\centering
		\caption{Performance comparison on the CFD dataset (detection metrics)}
		\begin{tabular}{ccccc}
			\toprule
			Method & Precision(\%) & Recall(\%) & mAP(\%) & Params(M) \\
			\midrule
			U-Net* & 82.5$\pm$1.3 & 80.8$\pm$1.8 & 72.4$\pm$1.5 & 31.0 \\
			Attention U-Net* & 77.7$\pm$1.5 & 84.2$\pm$1.4 & 71.0$\pm$1.6 & 32.5 \\
			SegNet* & 70.5$\pm$2.0 & 74.4$\pm$1.7 & 44.7$\pm$2.3 & 27.4 \\
			PSPNet* & 72.5$\pm$1.6 & 74.6$\pm$1.5 & 44.7$\pm$2.1 & 46.7 \\
			DeepCrack* & 65.5$\pm$2.1 & 76.0$\pm$1.9 & 76.6$\pm$1.4 & 38.5 \\
			BMS-YOLO & 89.7$\pm$0.5 & 87.4$\pm$0.7 & 88.1$\pm$0.4 & 3.02 \\
			\bottomrule
		\end{tabular}
		\label{tab:compare_cfd}
	\end{table*}

	Note: Models marked with * are segmentation models; their mAP values were computed by converting pixel-level predictions to bounding boxes using the minimum bounding rectangle strategy.
	
	\begin{table*}[htbp]
		\centering
		\caption{Performance comparison on the DeepCrack dataset (detection metrics)}
		\begin{tabular}{ccccc}
			\toprule
			Method & Precision(\%) & Recall(\%) & mAP(\%) & Params(M) \\
			\midrule
			U-Net* & 84.9$\pm$1.0 & 93.1$\pm$0.8 & 81.3$\pm$1.2 & 31.0 \\
			Attention U-Net* & 92.1$\pm$0.6 & 85.2$\pm$0.9 & 80.8$\pm$1.1 & 32.5 \\
			PSPNet* & 86.9$\pm$0.8 & 95.2$\pm$0.5 & 84.1$\pm$0.7 & 46.7 \\
			TCDNet* & 94.3$\pm$0.4 & 82.7$\pm$0.7 & 89.2$\pm$0.5 & 35.8 \\
			Swin-T* & 91.0$\pm$0.7 & 68.7$\pm$1.5 & 81.8$\pm$1.3 & 28.3 \\
			DeepCrack* & 91.9$\pm$0.5 & 78.3$\pm$1.0 & 86.0$\pm$0.8 & 38.5 \\
			BMS-YOLO & 88.1$\pm$0.6 & 86.8$\pm$0.8 & 82.4$\pm$0.5 & 3.02 \\
			\bottomrule
		\end{tabular}
		\label{tab:compare_deepcrack}
	\end{table*}

	Note: Models marked with * are segmentation models; their mAP values were computed by converting pixel-level predictions to bounding boxes using the minimum bounding rectangle strategy.
	
	On the CFD dataset, BMS-YOLO achieves the best precision (89.7\%) and recall (87.4\%), demonstrating good generalization ability. On DeepCrack, BMS-YOLO achieves competitive precision (88.1\%) and recall (86.8\%), though its mAP is lower than some specialized segmentation models (e.g., TCDNet 89.2\%, PSPNet 84.1\%). This is mainly attributed to dataset characteristics: DeepCrack contains more complex mesh-like cracks and low-contrast scenes, where semantic segmentation networks have an advantage due to their dense pixel-wise prediction capability. Nevertheless, as a box-supervised detection model with only 3.02M parameters, BMS-YOLO maintains balanced precision and recall, and its inference speed far exceeds all segmentation networks (typically <50 FPS). These results indicate that our method has good cross-dataset generalization while maintaining efficiency, though there is still room for improvement in extremely dense and complex texture scenarios.
	
	\subsection{Visual Analysis of Results}
	To intuitively demonstrate the detection performance of BMS-YOLO, Figure 1 shows typical detection results on different datasets, with qualitative comparisons against U-Net and PSPNet segmentation results.
	
	On the self-built dataset, BMS-YOLO accurately detects cracks of various orientations (horizontal, vertical, mesh-like), with localization accuracy significantly better than U-Net and PSPNet, which suffer from missing edges and background misclassification. For region-type distresses such as potholes and patches, BMS-YOLO also maintains high detection accuracy, with bounding boxes tightly adhering to distress boundaries, indicating effective activation of the regional expert in MorphSparseMoE. On the CFD dataset, under challenging conditions of complex pavement texture and uneven illumination, BMS-YOLO still detects the main crack body completely, while comparison methods produce fragmented false detections. On DeepCrack, facing dense mesh-like cracks, BMS-YOLO, although not achieving pixel-wise coverage like segmentation networks, provides effective global localization in the form of bounding boxes, and after topology fusion, fragmentation is significantly reduced. Overall, PSPNet is the most unstable across all scenarios, followed by U-Net; BMS-YOLO demonstrates the best detection completeness and localization accuracy in all scenarios.
	
	\subsection{Failure Case Analysis}
	Despite the overall excellent performance, BMS-YOLO still fails in some scenarios, as shown in Figure 2. Typical failure modes can be categorized into three types:
	
	(1) Extremely low-contrast cracks: When the grayscale difference between cracks and background pavement is very small ($\Delta I < 10$) and without significant texture variation, FDDE's high-frequency enhancement can partially improve responses but is insufficient to ensure stable detection, leading to misses. This problem is particularly prominent in shadow-covered areas, suggesting future work could combine illumination normalization preprocessing or introduce low-frequency compensation mechanisms.
	
	(2) Over-fragmentation in dense mesh-like cracks: In regions with dense mesh-like cracks, topology fusion can aggregate some fragments, but when crack spacing is extremely small and directions frequently interweave, the grouping thresholds of the Union-Find algorithm cannot simultaneously satisfy connectivity and distinctiveness, causing some boxes to remain fragmented or become over-fused. This indicates that the current directional consistency and distance criteria are not yet robust enough for extremely complex topologies.
	
	(3) Missed detection of extremely small crack fragments: For tiny crack segments with pixel width only 1-2 pixels and length less than 10 pixels, feature responses become extremely weak after multiple downsampling operations, and detection confidence falls below the threshold, resulting in misses. This is a common difficulty in small object detection; future work could combine higher-resolution branches in the feature pyramid or introduce super-resolution auxiliary modules to alleviate this issue.
	
	These failure cases point to directions for future improvement, including enhancing feature responses under low contrast, introducing finer topology modeling strategies, and strengthening feature representation for extremely small targets.
	
	\subsection{Full-Image Inference Results}
	To verify the model's deployment applicability in real scenes, we run the BMS-YOLO inference pipeline on original uncropped aerial images, with results shown in Figure 3. The model maintains stable detection performance on full images of 1024x1024 and above, providing reasonable localization for multi-scale distresses (from elongated cracks to large patches). The topology fusion post-processing effectively reduces redundant boxes for the same crack, making detection results more concise and reliable. Inference speed reaches 78 FPS (full image) without optimization. Based on preliminary INT8 quantization experiments on a subset of 50 images, we observed a 1.2\% mAP degradation (from 89.7\% to 88.5\%) with an estimated 2.5x speedup, suggesting that TensorRT-optimized deployment could achieve 300+ FPS on edge hardware such as NVIDIA Jetson Orin Nano. This demonstrates potential for real-time operation on UAV-borne edge devices.
	
	\section{Conclusion}
	This paper addresses the core issues of high-frequency detail loss, insufficient morphological adaptability, computational redundancy, and loss function unfriendliness to crack characteristics in existing YOLO models for pavement distress detection, and proposes a lightweight box-supervised detection model BMS-YOLO (Box-supervised Morphology-aware Sparse YOLO). Through systematic theoretical analysis and extensive experimental validation, the main conclusions are as follows:
	
	(1) The Frequency-Direction Detail Enhancement module effectively improves perception of subtle cracks. By decoupling the feature map into high- and low-frequency branches, introducing horizontal and vertical directional depthwise convolutions in the high-frequency branch, and dynamically fusing the two branches via an adaptive gating mechanism, the model significantly strengthens responses to fine crack edges with almost no additional computation. Ablation experiments show that the FDDE module independently contributes a 2.8\% mAP improvement, verifying the key role of explicitly enhancing high-frequency directional information for crack detection.
	
	(2) The Morphology-aware Sparse Mixture of Experts achieves morphology-adaptive modeling with low computational cost. By embedding four types of experts (horizontal-line, vertical-line, isotropic, and regional) inside the C2f module and introducing a lightweight routing network for dynamic top-k sparse activation, the model can selectively activate the most relevant expert branches according to local morphological patterns. Experiments show that MorphSparseMoE brings a 1.4\% mAP gain with only a 0.26M parameter increase, achieving an ideal balance between accuracy and efficiency. This design effectively solves the problem that homogeneous convolutions cannot simultaneously adapt to the extreme morphological differences between line-shaped cracks and region-shaped potholes.
	
	(3) ECA and the lightweight SPPF-CSPC structure synergistically achieve model compression and receptive field expansion. Channel attention is implemented via 1D convolution with adaptive kernel size, avoiding the information loss caused by dimensionality reduction in traditional SENet; the cross-stage partial connection and sequential parameter-free pooling design significantly reduce parameters (-17.7\%) and computation (-18.0\%) while enlarging the receptive field and enhancing multi-scale feature aggregation. The total number of parameters is only 3.02M, the smallest among all compared methods, indicating good potential for edge deployment.
	
	(4) The joint optimization of WIoU and box morphology consistency loss significantly improves supervision accuracy for elongated targets. WIoU alleviates gradient vanishing for small and non-overlapping boxes via center-distance-based dynamic focusing weights; the morphology consistency loss constrains aspect ratio and area in log space, with class-aware weighting that imposes higher morphology loss weights on crack classes. Ablation experiments show that WIoU alone contributes a 3.2\% mAP improvement, and the morphology loss further adds 0.7\%. Their combined use significantly enhances the model's ability to localize and preserve the shape of elongated cracks.
	
	(5) The confidence calibration and topology-guided box fusion post-processing strategy effectively solves the fragmentation problem in crack detection. Confidence calibration based on local edge density and shape priors corrects confidence biases in low-contrast scenes; topology fusion based on the Union-Find algorithm aggregates fragmented crack boxes using a low IoU threshold (0.12) and directional consistency, reducing the average number of duplicate boxes from 2.8 to 0.9 and improving recall by 3.5\%. This post-processing pipeline effectively improves the completeness and reliability of detection results without changing the forward computation of the model.
	
	(6) BMS-YOLO outperforms state-of-the-art methods in both accuracy and efficiency, with good generalization ability. On the self-built pavement distress dataset, the model achieves 89.7\% mAP and 136.8 FPS inference speed on a Quadro P2000 GPU, improving mAP by 6.4\% compared with baseline YOLOv8n while maintaining competitive inference speed. The model significantly outperforms detection-based baselines (YOLOv8n, YOLOv8n-Seg) in precision, recall, and mAP. Cross-domain validation on public CFD and DeepCrack datasets further demonstrates that BMS-YOLO exhibits good generalization capability while maintaining high efficiency, verifying the universality and robustness of each improved module design.
	
	In summary, the proposed BMS-YOLO model, through the systematic collaborative design of frequency-direction detail enhancement, morphology-aware sparse experts, lightweight attention and pooling structures, jointly optimized loss functions, and topology-aware post-processing, achieves efficient and accurate detection of pavement cracks and potholes under pure box supervision, providing a lightweight and highly reliable technical solution for real-time pavement distress inspection in practical engineering.
	
	\paragraph{Limitations and Future Work}
	Although BMS-YOLO achieves superior overall performance, there are still misses and fragmentation issues in detecting extremely low-contrast cracks (grayscale difference $\Delta I < 10$), dense mesh-like cracks, and extremely small cracks with size less than 10 pixels. In addition, the model's performance on high-complexity scenes like DeepCrack still has room for improvement. Future work will focus on the following directions: (1) introducing illumination normalization preprocessing or Retinex-based low-frequency compensation mechanisms to enhance crack responses in low-contrast scenes; (2) incorporating super-resolution auxiliary modules or attention-guided feature refinement in the high-resolution branches of the feature pyramid to strengthen feature representation for extremely small targets; (3) exploring graph neural networks or Transformer-based global context modeling to improve topology modeling in dense mesh-like crack scenes; (4) conducting model quantization (INT8) and structured pruning to further compress model size for efficient deployment on UAV-borne edge devices; and (5) exploring semi-supervised and weakly supervised learning paradigms to reduce reliance on large-scale accurate bounding box annotations and improve applicability in annotation-scarce scenarios.

	\paragraph{Scope and Applicability}
	It is important to clarify the intended application scope of BMS-YOLO. The model is designed for rapid screening and large-scale inspection of pavement distress --- i.e., detecting the presence and approximate location of cracks and potholes across extensive road networks. It is not intended for safety-critical structural assessment, where pixel-level precision (e.g., crack width measurement, crack length continuity analysis) is required. For such applications, box-level outputs could be augmented with a lightweight segmentation head or used to trigger follow-up high-resolution imaging at detected locations. This distinction is critical: bounding-box supervision offers a practical trade-off between annotation cost and detection accuracy for large-scale inspection, while segmentation-level supervision remains preferable for detailed structural analysis.
	
	% 参考文献（手动编写）
	\begin{thebibliography}{99}
		
		\bibitem{Adarsh2020}
		Adarsh, P., Rathi, P., Kumar, M.: YOLOv3-Tiny: Object Detection and Recognition using one stage improved model. In: 2020 6th International Conference on Advanced Computing and Communication Systems (ICACCS), pp. 687–694. IEEE (2020)
		
		\bibitem{Ali2021}
		Ali, L., Alnajjar, F., Jassmi, H.A., Gochoo, M., Almansoori, S., Khaled, F.A.: Performance evaluation of deep CNN-based crack detection and localization techniques for concrete structures. Sensors 21(5), 1688 (2021)
		
		\bibitem{Bochkovskiy2020}
		Bochkovskiy, A., Wang, C.-Y., Liao, H.-Y.M.: YOLOv4: Optimal speed and accuracy of object detection. arXiv preprint arXiv:2004.10934 (2020)
		
		\bibitem{Chen2017}
		Chen, L.-C., Papandreou, G., Schroff, F., Adam, H.: Rethinking atrous convolution for semantic image segmentation. arXiv preprint arXiv:1706.05587 (2017)
		
		\bibitem{Chollet2017}
		Chollet, F.: Xception: Deep learning with depthwise separable convolutions. In: Proceedings of the IEEE Conference on Computer Vision and Pattern Recognition, pp. 1251–1258 (2017)
		
		\bibitem{Dai2017}
		Dai, J., Qi, H., Xiong, Y., Li, Y., Zhang, G., Hu, X., Wei, Y.: Deformable convolutional networks. In: Proceedings of the IEEE International Conference on Computer Vision, pp. 764–773 (2017)
		
		\bibitem{Deng2020}
		Deng, J., Lu, Y., Lee, V.C.S.: Concrete crack detection with handwriting script interferences using faster region-based convolutional neural network. Computer-Aided Civil and Infrastructure Engineering 35(4), 373–388 (2020)
		
		\bibitem{Dorafshan2018}
		Dorafshan, S., Thomas, R.J., Maguire, M.: Comparison of deep convolutional neural networks and edge detectors for image-based crack detection in concrete. Construction and Building Materials 186, 1031–1045 (2018)
		
		\bibitem{Fan2024}
		Fan, Y., Li, H., Bao, Y.: Cycle-consistency-constrained few-shot learning framework for universal multi-type structural damage segmentation. Structural Health Monitoring (2024)
		
		\bibitem{Farhadi2018}
		Farhadi, A., Redmon, J.: YOLOv3: An incremental improvement. arXiv preprint arXiv:1804.02767 (2018)
		
		\bibitem{He2016}
		He, K., Zhang, X., Ren, S., Sun, J.: Deep residual learning for image recognition. In: Proceedings of the IEEE Conference on Computer Vision and Pattern Recognition, pp. 770–778 (2016)
		
		\bibitem{Howard2017}
		Howard, A.G., Zhu, M., Chen, B., Kalenichenko, D., Wang, W., Weyand, T., Andreetto, M., Adam, H.: MobileNets: Efficient convolutional neural networks for mobile vision applications. arXiv preprint arXiv:1704.04861 (2017)
		
		\bibitem{Hu2018}
		Hu, J., Shen, L., Sun, G.: Squeeze-and-excitation networks. In: Proceedings of the IEEE Conference on Computer Vision and Pattern Recognition, pp. 7132–7141 (2018)
		
		\bibitem{LeCun1998}
		LeCun, Y., Bottou, L., Bengio, Y., Haffner, P.: Gradient-based learning applied to document recognition. Proceedings of the IEEE 86(11), 2278–2324 (1998)
		
		\bibitem{Li2021}
		Li, D., Xie, Q., Gong, X., Zou, Q.: Automatic defect detection of metro tunnel surfaces using a vision-based inspection system. Advanced Engineering Informatics 47, 101206 (2021)
		
		\bibitem{Liu2021}
		Liu, Z., Lin, Y., Cao, Y., Hu, H., Wei, Y., Zhang, Z., Lin, S., Guo, B.: Swin transformer: Hierarchical vision transformer using shifted windows. In: Proceedings of the IEEE/CVF International Conference on Computer Vision, pp. 10012–10022 (2021)
		
		\bibitem{Redmon2016}
		Redmon, J., Divvala, S., Girshick, R., Farhadi, A.: You only look once: Unified, real-time object detection. In: Proceedings of the IEEE Conference on Computer Vision and Pattern Recognition, pp. 779–788 (2016)
		
		\bibitem{Ren2016}
		Ren, S., He, K., Girshick, R., Sun, J.: Faster R-CNN: Towards real-time object detection with region proposal networks. IEEE Transactions on Pattern Analysis and Machine Intelligence 39(6), 1137–1149 (2016)
		
		\bibitem{Rezatofighi2019}
		Rezatofighi, H., Tsoi, N., Gwak, J.Y., Sadeghian, A., Reid, I., Savarese, S.: Generalized intersection over union: A metric and a loss for bounding box regression. In: Proceedings of the IEEE/CVF Conference on Computer Vision and Pattern Recognition, pp. 658–666 (2019)
		
		\bibitem{Riquelme2021}
		Riquelme, C., Puigcerver, J., Mustafa, B., Neumann, M., Alayrac, J.-B., Zhai, X., Goh, A., Brock, A., Smith, S., Hechtman, K., et al.: Scaling vision with sparse mixture of experts. In: Advances in Neural Information Processing Systems, vol. 34, pp. 8583–8595 (2021)
		
		\bibitem{Ronneberger2015}
		Ronneberger, O., Fischer, P., Brox, T.: U-Net: Convolutional networks for biomedical image segmentation. In: International Conference on Medical Image Computing and Computer-Assisted Intervention, pp. 234–241. Springer (2015)
		
		\bibitem{Shazeer2017}
		Shazeer, N., Mirhoseini, A., Maziarz, K., Davis, A., Le, Q., Hinton, G., Dean, J.: Outrageously large neural networks: The sparsely-gated mixture-of-experts layer. arXiv preprint arXiv:1701.06538 (2017)
		
		\bibitem{Tan2019}
		Tan, M., Le, Q.: EfficientNet: Rethinking model scaling for convolutional neural networks. In: International Conference on Machine Learning, pp. 6105–6114. PMLR (2019)
		
		\bibitem{Tang2019}
		Tang, J., Mao, Y., Wang, J.: Multi-task enhanced dam crack image detection based on faster R-CNN. In: 2019 IEEE 4th International Conference on Image, Vision and Computing (ICIVC), pp. 336–340. IEEE (2019)
		
		\bibitem{Jocher2023}
		Jocher, G., Chaurasia, A., Qiu, J.: YOLOv8: Ultralytics YOLO object detection, model training, and deployment tools. https://github.com/ultralytics/ultralytics (2023)

		\bibitem{Sui2022}
		Sui, X., Wang, J., Li, H., Zhou, R., Xu, Y., Jin, X.: Wise-IoU: Thinking the bias of IoU for robust object detection. arXiv preprint arXiv:2208.10791 (2022)

		\bibitem{Zhou2018}
		Zhou, X., Zhuo, J.: Accurate crack detection via adaptive context representation. In: Proceedings of the IEEE/CVF Conference on Computer Vision and Pattern Recognition Workshops, pp. 2303–2304 (2018)

		\bibitem{Li2020}
		Li, X., Wang, W., Hu, X.: Generalized intersection over union: A diversity-preserving IoU loss. IEEE Transactions on Image Processing 29, 5678–5690 (2020)

		\bibitem{Murillo2013}
		Murillo, C., Tapia, J., Guerrero, J., Maldonado, R., Sandoval, R.: Survey on UAV-based crack detection for infrastructure inspection. Automation in Construction 156, 105128 (2024)

		\bibitem{Wang2024}
		Wang, Y., Zhang, L., Li, R.: YOLO-based pavement distress detection: A comprehensive review. Construction and Building Materials 412, 134256 (2024)

		\bibitem{Varghese2024}
		Varghese, R., Sambath, M.: YOLOv8: A novel object detection algorithm with enhanced performance and robustness. In: International Conference on Advances in Data Engineering and Intelligent Computing Systems, pp. 1–6. IEEE (2024)
		
		\bibitem{Wang2021}
		Wang, J., Xu, C., Yang, W., Yu, L.: A normalized Gaussian Wasserstein distance for tiny object detection. arXiv preprint arXiv:2110.13389 (2021)
		
		\bibitem{Wang2020}
		Wang, Q., Wu, B., Zhu, P., Li, P., Zuo, W., Hu, Q.: ECA-Net: Efficient channel attention for deep convolutional neural networks. In: Proceedings of the IEEE/CVF Conference on Computer Vision and Pattern Recognition, pp. 11534–11542 (2020)
		
		\bibitem{Wang2023}
		Wang, C.-Y., Bochkovskiy, A., Liao, H.-Y.M.: YOLOv7: Trainable bag-of-freebies sets new state-of-the-art for real-time object detectors. In: Proceedings of the IEEE/CVF Conference on Computer Vision and Pattern Recognition, pp. 7464–7473 (2023)
		
		\bibitem{Xie2017}
		Xie, S., Girshick, R., Dollár, P., Tu, Z., He, K.: Aggregated residual transformations for deep neural networks. In: Proceedings of the IEEE Conference on Computer Vision and Pattern Recognition, pp. 1492–1500 (2017)
		
		\bibitem{Xu2023}
		Xu, Y., Fan, Y., Li, H.: Lightweight semantic segmentation of complex structural damage recognition for actual bridges. Structural Health Monitoring 22(5), 3250–3269 (2023)
		
		\bibitem{Yu2015}
		Yu, F., Koltun, V.: Multi-scale context aggregation by dilated convolutions. arXiv preprint arXiv:1511.07122 (2015)
		
		\bibitem{Yu2016}
		Yu, J., Jiang, Y., Wang, Z., Li, Z.: UnitBox: An advanced object detection network. In: Proceedings of the 24th ACM International Conference on Multimedia, pp. 516–520 (2016)
		
		\bibitem{Zhang2018}
		Zhang, X., Zhou, X., Lin, M., Sun, J.: ShuffleNet: An extremely efficient convolutional neural network for mobile devices. In: Proceedings of the IEEE Conference on Computer Vision and Pattern Recognition, pp. 6848–6856 (2018)
		
		\bibitem{Zhang2024}
		Zhang, Y., Liu, C.: Crack segmentation using discrete cosine transform in shadow environments. Automation in Construction 166, 105646 (2024)
		
		\bibitem{Zhao2017}
		Zhao, H., Shi, J., Qi, X., Wang, X., Jia, J.: Pyramid scene parsing network. In: Proceedings of the IEEE Conference on Computer Vision and Pattern Recognition, pp. 2881–2890 (2017)
		
		\bibitem{Zheng2020}
		Zheng, Z., Wang, P., Liu, W., Li, J., Ye, R., Ren, D.: Distance-IoU loss: Faster and better learning for bounding box regression. In: Proceedings of the AAAI Conference on Artificial Intelligence, vol. 34(7), pp. 12993–13000 (2020)
		
		\bibitem{Zou2018}
		Zou, Q., Zhang, Z., Li, Q., Wang, H.: DeepCrack: Learning hierarchical convolutional features for crack detection. IEEE Transactions on Image Processing 28, 1498–1512 (2018)
		
	\end{thebibliography}
	
\end{document}